\section{Diseño de Pruebas de Validación}

La validación del diseño y despliegue requiere comprobar que la topología
física se ajusta al plan lógico, que los servicios son accesibles desde todos
los segmentos, y que los mecanismos de seguridad operan conforme a lo especificado.

Este apartado describe el conjunto de pruebas, sus procedimientos, criterios de
éxito y herramientas utilizadas. Las pruebas se organizan por categoría
funcional: direccionamiento, enrutamiento, servicios, y seguridad.

\subsection{Matriz de Pruebas de Direccionamiento}

\subsubsection{Prueba 1.1: Validación de Asignación IPv4}

\textbf{Objetivo:} Verificar que cada interfaz de router y servidor recibió la
dirección IPv4 estática correcta según el plan de direccionamiento.

\textbf{Procedimiento:}
\begin{enumerate}
\item Acceder al router de la Sede B
\item Ejecutar comando \texttt{show ip interface brief} para listar todas las interfaces
\item Verificar que cada interfaz coincida con la asignación en tabla~\ref{tab:ipv4}
\item En los servidores Linux, ejecutar \texttt{ip addr show} y validar direcciones IPv4
\end{enumerate}

\textbf{Criterio de Éxito:}
\begin{itemize}
\item Todas las interfaces tienen dirección IPv4 del rango correcto
\item Máscara de red (prefijo) coincide con el plan de VLSM
\end{itemize}


\subsubsection{Prueba 1.2: Validación de Asignación IPv6}

\textbf{Objetivo:} Verificar que cada interfaz de router recibió la dirección IPv6
estática correcta y que SLAAC funciona en los hosts clientes.

\textbf{Procedimiento:}
\begin{enumerate}
\item En router, ejecutar \texttt{show ipv6 interface brief}
\item Validar que cada interfaz tiene dirección IPv6 del prefijo correcto (tabla~\ref{tab:ipv6})
\item En cliente Windows, abrir PowerShell y ejecutar \texttt{ipconfig /all}
\item Verificar que el cliente tiene dirección IPv6 que comienza con el prefijo anunciado
\end{enumerate}

\textbf{Criterio de Éxito:}
\begin{itemize}
\item Router tiene dirección IPv6 \texttt{/64} del rango correcto en cada interfaz
\item Clientes reciben dirección IPv6 con prefijo correcto mediante SLAAC
\end{itemize}


\subsubsection{Prueba 1.3: DHCP en IPv4}

\textbf{Objetivo:} Verificar que el servidor DHCP asigna direcciones dentro del rango
correcto de cada LAN y que los clientes las adquieren automáticamente.

\textbf{Procedimiento:}
\begin{enumerate}
\item En un cliente de cada LAN, ejecutar \texttt{ipconfig /release} y luego \texttt{ipconfig /renew}
\item Capturar la dirección asignada con \texttt{ipconfig}
\item Verificar que está dentro del rango de usuario final para esa LAN:
  \begin{itemize}
  \item LAN 1: \ipRangoLanUno
  \item LAN 2: \ipRangoLanDos
  \item LAN 3: \ipRangoLanTres
  \end{itemize}
\item Ejecutar \texttt{ipconfig /all} para verificar gateway y máscara
\end{enumerate}

\textbf{Criterio de Éxito:}
\begin{itemize}
\item Dirección asignada está dentro del rango de subred de usuario final
\item Gateway coincide con la dirección del router en esa LAN
\item Máscara de red es correcta (255.255.252.0 para LAN 1/2, 255.255.255.0 para LAN 3)
\end{itemize}


\subsection{Matriz de Pruebas de Conectividad y Enrutamiento}

\subsubsection{Prueba 2.1: Ping entre Hosts en la Misma LAN}

\textbf{Objetivo:} Verificar conectividad L2 dentro de cada segmento.

\textbf{Procedimiento:}
\begin{enumerate}
\item Desde Host en LAN 1, ejecutar \texttt{ping} a otro Host en LAN 1
\item Capturar tiempo de respuesta y pérdida de paquetes
\item Repetir para LAN 2 y LAN 3
\end{enumerate}

\textbf{Criterio de Éxito:}
\begin{itemize}
\item 100\% de paquetes respondidos
\item Latencia menor a 10 ms (esperable en LAN local)
\end{itemize}

\subsubsection{Prueba 2.2: Ping entre Hosts en LANs Diferentes}

\textbf{Objetivo:} Verificar que el enrutamiento IPv4 funciona y que los paquetes
traversan correctamente entre sedes.

\textbf{Procedimiento:}
\begin{enumerate}
\item Desde Host en LAN 1 (ej: \ipHostLanUno), ejecutar \texttt{ping \ipHostLanDos}
\item Capturar respuesta y latencia
\item Desde LAN 2, hacer ping a LAN 3
\item Desde LAN 3, hacer ping a LAN 1
\end{enumerate}

\textbf{Criterio de Éxito:}
\begin{itemize}
\item Todos los pings reciben respuesta
\item Latencia entre 5--50 ms (según topología física real)
\end{itemize}


\subsubsection{Prueba 2.3: Ping IPv6 entre Hosts}

\textbf{Objetivo:} Validar que IPv6 routing funciona.

\textbf{Procedimiento:}
\begin{enumerate}
\item Desde cliente en LAN 1, ejecutar \texttt{ping} a dirección IPv6 de cliente en LAN 2
\item Ej: \texttt{ping \ipHostIpvSeisLanDos}
\item Repetir entre otros pares de LANs
\end{enumerate}

\textbf{Criterio de Éxito:}
\begin{itemize}
\item IPv6 ping recibe respuesta
\item Latencia similar a IPv4
\item Ninguna pérdida de paquetes
\end{itemize}


\subsection{Matriz de Pruebas de Servicios}

\subsubsection{Prueba 3.1: Accesibilidad del Servidor Web}

\textbf{Objetivo:} Verificar que el servidor web (Nginx) está operativo y accesible
desde todos los segmentos en IPv4 e IPv6.

\textbf{Procedimiento:}
\begin{enumerate}
\item Desde cliente en LAN 1, abrir navegador y acceder a \texttt{http://\ipServidorWeb}
\item Capturar página de bienvenida
\item Repetir desde clientes en LAN 2 y LAN 3
\item Abrir navegador y acceder a \texttt{http://[\ipServidorWebIpvSeis]}
\item Ejecutar \texttt{curl}: \texttt{curl -v http://\ipServidorWeb}
\end{enumerate}

\textbf{Criterio de Éxito:}
\begin{itemize}
\item Página carga correctamente en todos los clientes
\item Código HTTP 200 OK en respuesta
\item Ambas direcciones IPv4 e IPv6 responden
\item Latencia menor a 500 ms
\end{itemize}


\subsubsection{Prueba 3.2: Conectividad del Servidor de Correo}

\textbf{Objetivo:} Verificar que los servicios de correo (SMTP e IMAP) están disponibles.

\textbf{Procedimiento:}
\begin{enumerate}
\item Enviar un correo de prueba a otro usuario desde el portal
\item Verificar que el correo aparece en la bandeja de entrada del destinatario
\end{enumerate}

\textbf{Criterio de Éxito:}
\begin{itemize}
\item El portal web puede autenticarse y listar la bandeja de entrada vía IMAP
\item El correo de prueba se entrega y es visible en el portal del destinatario
\end{itemize}


\subsubsection{Prueba 3.3: Servidor de Chat}

\textbf{Objetivo:} Verificar que el servidor de chat está operativo y acepta
conexiones de múltiples clientes simultáneos.

\textbf{Procedimiento:}
\begin{enumerate}
\item En servidor Linux, verificar que el demonio de chat está corriendo:
  \texttt{ps aux | grep chat}
\item Verificar que está escuchando en el puerto asignado:
  \texttt{ss -tuln | grep 8080}
\item En cliente Windows 1, conectar al servidor: ejecutar cliente con
  \texttt{./chat-client \ipServidorChat\xspace 8080}
\item En cliente Windows 2, conectar al mismo servidor
\item En cliente 1, ejecutar \texttt{/list} y verificar que cliente 2 aparece
\item Enviar mensaje desde cliente 1 y verificar que cliente 2 lo recibe
\item Validar comando \texttt{/save} para guardar historial en archivo
\item Conectar cliente desde LAN 2 y LAN 3 y repetir pruebas
\end{enumerate}

\textbf{Criterio de Éxito:}
\begin{itemize}
\item Servidor está en estado \texttt{LISTEN} en puerto correcto
\item Múltiples clientes se conectan simultáneamente
\item Comando \texttt{/list} muestra todos los usuarios conectados
\item Mensajes se entregan sin retraso entre clientes
\item Historial se exporta correctamente a archivo
\item Funciona desde cualquier LAN (pruebas de conectividad cruzada)
\end{itemize}


\subsection{Matriz de Pruebas de Seguridad}

\subsubsection{Prueba 4.1: SSH Disponible, Telnet Bloqueado}

\textbf{Objetivo:} Verificar que SSHv2 está activo y Telnet deshabilitado en routers.

\textbf{Procedimiento:}
\begin{enumerate}
\item Desde cliente, intentar \texttt{telnet \ipRouterUno 23}
\item Verificar que \textbf{no} responde (conexión rechazada o timeout)
\item Ejecutar \texttt{ssh admin@\ipRouterUno}
\item Ingresar contraseña
\item Verificar que SSH conecta exitosamente
\item Ejecutar \texttt{show version} en router para confirmar acceso
\end{enumerate}

\textbf{Criterio de Éxito:}
\begin{itemize}
\item Telnet no responde (puerto 23 cerrado)
\item SSH responde y permite ingreso
\item Sesión SSH cifra toda la comunicación
\end{itemize}

\textbf{Herramientas:} \texttt{telnet}, \texttt{ssh}, cliente terminal

\subsubsection{Prueba 4.2: Firewall en Servidor Linux}

\textbf{Objetivo:} Verificar que solo los puertos requeridos están abiertos.

\textbf{Procedimiento:}
\begin{enumerate}
\item Ejecutar \texttt{sudo ufw status} en los servidores Linux
\item Verificar que solo están abiertos: 2222 (SSH), 25 (SMTP),
  143 (IMAP), y el puerto 8080 del servidor de chat. 
  Para el servidor web debe de estar abierto el puerto 80 (HTTP)
\end{enumerate}

\textbf{Criterio de Éxito:}
\begin{itemize}
\item Firewall está activo y en modo \texttt{active}
\item Solo puertos permitidos responden
\end{itemize}


\subsubsection{Prueba 4.3: Contraseñas Cifradas en Routers}

\textbf{Objetivo:} Verificar que las contraseñas en los archivos de configuración
se almacenan en forma cifrada (no texto plano).

\textbf{Procedimiento:}
\begin{enumerate}
\item Entrar a la consola de un router (ej: Router Sede B)
\item Ejecutar \texttt{show run} para ver configuración
\item Buscar sección de contraseñas (enable secret, password)
\item Verificar que \textbf{enable secret} comienza con \texttt{\$5\$} o \texttt{\$6\$}
  (indicador de hash cifrado)
\item Verificar que las contraseñas en líneas VTY no aparecen en texto plano
\end{enumerate}

\textbf{Criterio de Éxito:}
\begin{itemize}
\item \texttt{enable secret} está cifrado (comienza con \$)
\item Contraseñas de VTY no están en texto plano en configuración
\item Usuario no puede descifrar las contraseñas desde archivo de respaldo
\end{itemize}


